\section{Results}
\label{sec:results}
Figure~\ref{fig:ops} contains the paper-level evidence. The companion repository mirrors this section and provides the full numerical tables, all confidence intervals and every sensitivity analysis.

\subsection{Recovery grows with requested duration}
Severity~2 provides the clearest test. At $3.84$\,s, recovery gain is $\Rs=+0.085$ \ci{+0.066}{+0.105}. At $10.24$\,s it is $\Rn=+0.244$ \ci{+0.215}{+0.273}. Their pre-specified difference is $J=+0.159$ \ci{+0.131}{+0.187}, so the fine-tuning stage contributes substantially more at the duration on which it was performed. The seam sensitivity gives the same conclusion.

The four-point sweep shows that this is not an endpoint accident. Recovery rises at every tested duration between $3.84$ and $10.24$\,s, and every adjacent increase is resolved under the primary scorer. The scale of the change is also meaningful. Relative to real audio, fine-tuning closes $33\%$ of P's gap at $3.84$\,s and $63\%$ at $10.24$\,s. Relative to the matched dense model, the corresponding ratios are $44\%$ and $82\%$. The intermediate values and their intervals are reported in the companion repository~\cite{bibbo2026repo}.

The result is not explained by a generic tendency for CLAP to rise with longer audio. Dense generations and real audio also change with duration, but P and \PFT{} change differently. Chance correction strengthens rather than removes the interaction. The published sampler reproduces it, and Human-CLAP plus two event-level measures outside the CLAP family show the same larger-native-than-short pattern. A Holm correction over the registered severity~2 family leaves the duration conclusions unchanged. Exact robustness contrasts are kept in the companion repository because their role here is corroboration rather than a second headline result~\cite{bibbo2026repo}.

Severity~1 is directionally consistent but weaker. Its raw interaction is $+0.044$ \ci{-0.001}{+0.087}, and the follow-up was powered only for a larger effect. Chance correction yields a positive interaction, but we base the duration claim on the replicated severity~2 evidence rather than on this borderline result.

\subsection{Recovery is specific to the prompt domain}
The domain comparison changes the interpretation of the duration result. Recovery is clearly positive on AudioCaps at both tested durations, but the held-out hip-hop battery shows no resolved alignment gain at either one. At $10.24$\,s, for example, AudioCaps recovery is $+0.244$ while music recovery is $+0.005$ \ci{-0.028}{+0.039}. The large native-duration gain is therefore not a general improvement that transfers to the held-out captions. Full short-duration values are reported in the companion repository~\cite{bibbo2026repo}.

The anchors make the contrast more intuitive. After fine-tuning, the music cells sit only about $0.02$ to $0.03$ CLAP units above their shuffled-caption floor, whereas severity~2 AudioCaps at $10.24$\,s sits about $0.32$ above chance. Severity~1 showed a negative music gain, but that penalty did not replicate at severity~2. The stronger and reproducible conclusion is the absence of recovery on the held-out domain.

Our small blind listening check also separates alignment from perceptual quality. At $10.24$\,s on AudioCaps, \PFT{} was preferred on 6 of 8 pairs and P mostly sounded like noise. At $3.84$\,s both systems sounded like noise. On music, \PFT{} sounded more musical even though neither system followed the long captions well. This disagreement is one reason we interpret the primary endpoint specifically as text-audio alignment rather than as overall audio quality.

\subsection{The short-duration deficit comes from generation}
A natural explanation for the duration effect is that \PFT{} might place useful content late in a $10.24$\,s clip. FineLAP does not support that account. At severity~2 the grounding gain is distributed across the clip, and the difference between early and late gain is $T=-0.002$ \ci{-0.024}{+0.020}. Recovery is therefore not simply back-loaded.

A crop test then asks whether the deficit could come from how CLAP scores short audio. The first $3.84$\,s of a native-duration generation yields recovery $R_{\mathrm{crop}}=+0.172$ \ci{+0.150}{+0.194}, which is $+0.087$ \ci{+0.065}{+0.110} larger than recovery from a separately generated $3.84$\,s clip. The operating-point effect therefore arises primarily from \emph{generating} at short duration rather than from scoring a short window.

